\chapter{Corneal Abrasion}

\section*{Definition}
A corneal abrasion is a defect of the corneal epithelium exposing the underlying Bowman's layer and stroma, typically caused by mechanical trauma; healing of uncomplicated abrasions normally occurs within 24--72 hours.

\section*{Pathophysiology}
Loss of epithelial cell integrity disrupts the corneal barrier, exposing corneal nerve endings (the most densely innervated tissue in the body), which produces severe pain and reflex blepharospasm. The epithelial defect triggers a healing cascade involving cellular migration, proliferation, and adhesion.

\section*{Causes}
\begin{itemize}
  \item \textbf{Mechanical trauma:} Fingernail scratch, tree branch, paper edge, foreign body
  \item \textbf{Contact lens related:} Overnight wear, poor hygiene, dry lens insertion
  \item \textbf{Chemical injury:} Alkali or acid splash leading to epithelial sloughing
  \item \textbf{Ultraviolet keratitis:} Prolonged UV exposure (welding arc, tanning bed, snow reflection)
  \item \textbf{Iatrogenic:} Corneal scraping or post-surgical epithelial defect
\end{itemize}

\section*{When to See a Doctor}
See your doctor if:
\begin{itemize}
  \item Severe eye pain, photophobia, and foreign body sensation after known trauma
  \item Suspected corneal abrasion in a contact lens wearer (high-risk for microbial keratitis)
  \item Worsening pain or discharge 24 hours after initial injury
  \item Visual acuity decline not explained by epithelial defect alone
  \item Presence of white infiltrate on cornea concerning for infection
\end{itemize}

\section*{Related Symptoms}
\begin{itemize}
  \item Foreign body sensation and tearing
  \item Photophobia (ocular)
  \item Blurred vision (transient due to tear film disruption)
  \item Blepharospasm (involuntary eye closure)
  \item Ciliary flush (circumcorneal injection)
\end{itemize}
\textbf{Important:} Contact lens-associated abrasions require mandatory topical antibiotic coverage for Pseudomonas aeruginosa, which can cause rapid, sight-threatening keratitis.

\section*{Self-Care / Home Management}
\begin{itemize}
  \item Do not rub the affected eye
  \item Artificial tears (preservative-free) for comfort
  \item Removal of contact lenses until complete healing confirmed
  \item Sunglasses for photophobia relief
  \item Avoid eye patching (shown to delay healing and increase infection risk)
\end{itemize}

\section*{How Your Doctor Will Evaluate You}
\begin{itemize}
  \item Fluorescein instillation and cobalt blue light examination to delineate epithelial defect
  \item Slit-lamp biomicroscopy to assess depth, infiltrate, and anterior chamber reaction
  \item Seidel test to rule out full-thickness corneal perforation
  \item Everted lid examination to exclude retained subtarsal foreign body
  \item Corneal culture if infiltrate, hypopyon, or contact lens use
\end{itemize}

\section*{Treatment Options}
\begin{itemize}
  \item Topical broad-spectrum antibiotic (fluoroquinolone or aminoglycoside) qid for 3--5 days
  \item Topical cycloplegic (homatropine or cyclopentolate) for pain relief from ciliary spasm
  \item Oral analgesics (acetaminophen or NSAIDs) for pain management
  \item Debridement of loose epithelium with a sterile cotton swab
  \item Bandage contact lens for large, painful abrasions (with concurrent antibiotic prophylaxis)
  \item Referral for corneal cyanoacrylate glue or amniotic membrane graft if persistent epithelial defect
\end{itemize}

\clearpage
